\clearpage
\section{Kiểm tra package và đánh giá kết quả}
\subsection{Kiểm tra mã nguồn ngày 20/10/2026}
\begin{table}[H]\centering\small
\caption{Phạm vi và kết quả kiểm tra trong lần rà soát hiện tại}\label{tab:checks}
\begin{tabularx}{\linewidth}{@{}P{4.3cm}Y@{}}\toprule
Hạng mục&Kết quả và phạm vi\\\midrule
Unit test&20/20 passed. Cấu hình/hash/vùng trống/quaternion, xác nhận bbox, geometry scene và shutdown navigation.\\
ROS timer/shutdown&Chạy test ở domain riêng 67; không gửi goal. Kiểm tra timeout khi clock mô phỏng đứng yên và giữ ưu tiên Ctrl+C.\\
Build&\code{colcon build --symlink-install --packages-select turtlebot4_lab05} thành công.\\
Launch/executable&Hai launch đọc được arguments; năm executable được đăng ký. Không khởi chạy simulator hay mission trong lần rà soát này.\\
Scene hiện tại&8/8 vật thể qua validator map/unknown/tường SDF/khoảng hở từng cặp. Hash map khớp dữ liệu đã lưu.\\
Bằng chứng runtime có sẵn&Đối chiếu JSON, UUID, ảnh RViz/YOLO, quỹ đạo và số liệu lịch sử; không gán thành benchmark phiên bản mới.\\\bottomrule
\end{tabularx}\end{table}
\begin{lstlisting}[language=bash,caption={Lệnh tái lập kiểm tra không điều khiển robot},label={lst:checks}]
cd ~/turtlebot4_lab_ws
source lab_env.bash
source src/turtlebot4_lab05/scripts/lab05_env.bash
ROS_DOMAIN_ID=67 python -m pytest src/turtlebot4_lab05/test -q
colcon build --symlink-install --packages-select turtlebot4_lab05
python src/turtlebot4_lab05/scripts/check_scene.py
ros2 launch turtlebot4_lab05 bringup.launch.py --show-args
ros2 launch turtlebot4_lab05 yolo.launch.py --show-args
ros2 pkg executables turtlebot4_lab05
\end{lstlisting}
Các test dùng box tổng hợp kiểm tra logic xác nhận, không thay thế một dataset ảnh ở nhiều góc. Test scene kiểm tra mô hình hình học đã khai báo, không thay thế chạy robot trên bố trí mới. Test Ctrl+C mô phỏng state của client; không phải thử cắt mạng Nav2 thật.

\subsection{Cách đối chiếu một mission thành công}
Trong log \code{scattered_mission_events.json}, ba result có UUID khác nhau, status 4 và thứ tự Goal\_1--Goal\_2--Goal\_3. Tóm tắt ghi completed=3 và yolo\_interrupted=false. Observer đồng thời thấy 615 inference lúc odometry vượt ngưỡng chuyển động. Kết hợp ba loại bằng chứng mạnh hơn việc nhìn riêng ảnh terminal hoặc riêng ảnh camera.

Khả năng hoàn thành ba goal đã được kiểm chứng ở scene lịch sử. Sau khi thay balô bằng xe đạp, thêm vật thể và nâng model/bộ lọc, cần một lượt runtime mới nếu muốn kết luận cùng chỉ số cho cấu hình hiện tại. Video liên kết ở báo cáo là tài nguyên người thực hiện cung cấp; các bảng số liệu tự động vẫn chỉ dẫn tới các tệp log tương ứng.
\clearpage

\section{Thảo luận, kiến thức rút ra và hướng phát triển}
\subsection{Những điểm cần hiểu khi giải thích chương trình}
\begin{enumerate}
\item \textbf{Localization khác mapping:} mapping tạo/cập nhật map; AMCL sử dụng map có sẵn để ước lượng pose. Demo này nạp PGM cũ và không chạy SLAM.
\item \textbf{Goal khác vận tốc:} node Python gửi mục tiêu mức cao qua NavigateToPose. Nav2 lập đường, chấm quỹ đạo, phát vận tốc; Create 3/bridge thực thi. Có ba goal không có nghĩa viết ba lệnh cmd\_vel cố định.
\item \textbf{Map frame khác odom frame:} map dùng định vị toàn cục; odom liên tục cục bộ. AMCL cung cấp hiệu chỉnh giữa chúng; pose/goal và TF phải nhất quán.
\item \textbf{Accepted khác succeeded:} server nhận goal chưa đồng nghĩa robot đã tới. Code chỉ tăng số mục tiêu hoàn thành khi action terminal status bằng 4.
\item \textbf{YOLO không thay costmap:} ảnh phát hiện người/ghế là thông tin thị giác 2D. LiDAR/costmap mới là đầu vào tránh vật cản trong triển khai này.
\item \textbf{Đồng thời không đồng nghĩa cùng tần số:} camera có thể phát nhanh hơn inference. Latest-frame worker ưu tiên độ mới, bỏ frame cũ để tránh tích lũy trễ.
\item \textbf{Confidence không phải accuracy:} một nhãn confidence cao vẫn có thể sai trong domain mô phỏng. Cần dữ liệu có ground truth và metric theo lớp để định lượng độ chính xác.
\end{enumerate}
\subsection{Giới hạn và cách cải thiện có thể kiểm chứng}
YOLO pretrained gặp domain gap khi mesh/texture, bóng và góc nhìn khác ảnh huấn luyện. Thay balô dễ nhầm bằng xe đạp làm thay đổi nội dung scene; không tự chứng minh mọi góc nhìn của xe đạp đều nhận đúng. Cần lấy nhiều frame ở các đoạn đường, gán nhãn, thống kê false positive/false negative rồi so sánh YOLO11n/11s trên cùng dữ liệu.

Bộ lọc hai frame giảm nhiễu thoáng qua nhưng có thể bỏ vật thể xuất hiện ngắn hoặc di chuyển nhanh trong ảnh. Hạ ngưỡng cây tăng recall tiềm năng đồng thời có thể tăng nhầm nhãn. Có thể tinh chỉnh từ tập validation hoặc fine-tune bằng ảnh mô phỏng đa góc, giữ tập đánh giá độc lập; không chọn vài confidence cực đại để báo cáo là độ chính xác tổng thể.

Muốn liên hệ vật thể với map cần depth hoặc hình học camera, TF camera--base--map, đồng bộ thời gian và ước lượng bất định. Chức năng đó chưa có trong bài hiện tại. Trên robot thật cũng cần xác minh lại sensor/topic, TF, tốc độ và đặc tính vật lý; kết quả maze mô phỏng không tự chuyển thành kiểm định thực địa.
\subsection{Kết luận}
Chương trình hoàn chỉnh kết nối map server--AMCL--Nav2 và một nhánh YOLO độc lập, có ba goal đo từ bài trước, quản lý trạng thái/action rõ ràng và lưu bằng chứng. Các log lịch sử xác nhận ba goal SUCCEEDED liên tiếp cùng inference khi robot chuyển động; bản hiện tại đã qua unit/build/geometry. Cấu trúc này cho phép tái lập demo, giải thích từng node/topic/action và tiếp tục đánh giá nhận diện có kiểm soát.
\clearpage

\section*{Tài liệu tham khảo và tài nguyên nộp bài}
\addcontentsline{toc}{section}{Tài liệu tham khảo và tài nguyên nộp bài}
\begin{enumerate}[label={[\arabic*]}]
\item Mã nguồn và cấu hình bài 05: \ghlink{}. Các nhận xét về state machine, thresholds, QoS, hash và xử lý ảnh được đối chiếu trực tiếp với source trong package.
\item Video do người thực hiện cung cấp: \vidlink{}. Liên kết ngắn trong PDF trỏ trực tiếp tới tệp Drive.
\item Clearpath Robotics, \href{https://turtlebot.github.io/turtlebot4-user-manual/tutorials/navigation.html}{TurtleBot 4 User Manual -- Navigation}. Tham khảo cấu trúc localization và navigation chính thức.
\item Nav2, \href{https://github.com/ros-navigation/navigation2/tree/humble/nav2_amcl}{AMCL Configuration}. Nguyên lý định vị và ý nghĩa tham số AMCL; giá trị trong báo cáo lấy từ bản Humble đang cài.
\item Nav2, \href{https://github.com/ros-navigation/navigation2/tree/humble/nav2_navfn_planner}{NavFn Planner Configuration} và \href{https://github.com/ros-navigation/navigation2/tree/humble/nav2_dwb_controller}{DWB Controller Configuration}. Đối chiếu plugin NavFn/DWB với cấu hình tại máy.
\item ROS 2 Humble, \code{nav2_msgs/action/NavigateToPose.action}, bản cài trong \code{/opt/ros/humble/share/nav2_msgs/action/}. Định nghĩa goal, feedback và trường result; status kết thúc do ROS action wrapper cung cấp.
\item Ultralytics, \href{https://docs.ultralytics.com/models/yolo11/}{YOLO11} và \href{https://docs.ultralytics.com/modes/predict/}{Predict mode}. Cấu hình model.predict và trường bbox của Results; phiên bản thư viện sử dụng là 8.3.221.
\item Bằng chứng nội bộ: \code{validation/demo_metrics.json}, \code{scattered_demo_summary.json}, \code{scattered_mission_events.json}, \code{scene_layout.json}; ảnh native trong thư mục validation. Các số liệu có cùng nguồn được mô tả trong chú thích bảng/hình.
\end{enumerate}
\note{Tái lập báo cáo: mã LaTeX và ảnh nằm trong \code{report/} của package. Có thể biên dịch bằng pdfLaTeX hai lần hoặc Tectonic. Các sơ đồ pipeline là TikZ vector; sơ đồ map và biểu đồ số liệu được dựng bằng script từ map/config/JSON. Ảnh RViz/camera giữ nguồn thực nghiệm và được chú thích phiên bản.}
